% ============================================================
% Formulario, argomento 20: test di ipotesi di Neyman-Pearson (dispensa, capitolo 20)
% ============================================================
\argomento{Test di ipotesi: il test di Neyman-Pearson}{Dispensa: capitolo 20}[Test di Neyman-Pearson]

\voce{errorispecie}{Errori di prima e di seconda specie}{$H_0$ ipotesi nulla, $H_1$ alternativa; in $\Omega_1$ si
  rifiuta $H_0$}
  {\alpha = P\{D(\vett{X}) = 1 \mid H_0\} \qquad \beta = P\{D(\vett{X}) = 0 \mid H_1\}}
  {\variante{falso allarme (prima specie), dati continui}{\alpha = \int_{\Omega_1} f_{\vett{X} \mid H_0}(\vett{x}
     \mid H_0) \dd\vett{x}}
   \variante{dati discreti}{\alpha = \sum_{\vett{x} \in \Omega_1} p_{\vett{X} \mid H_0}(\vett{x} \mid H_0)}
   \variante{potenza del test (corretta rivelazione)}{1 - \beta = P\{D(\vett{X}) = 1 \mid H_1\}}
   \variante{potenza, dati continui}{1 - \beta = \int_{\Omega_1} f_{\vett{X} \mid H_1}(\vett{x} \mid H_1)
     \dd\vett{x}}
   \variante{$\beta$ è la mancata rivelazione (seconda specie); allargando $\Omega_1$ crescono sia la potenza sia i
     falsi allarmi}{}}

\voce{neymanpearson}{Lemma di Neyman-Pearson}{il test più potente con falso allarme non superiore ad $\alpha$}
  {L(\vett{x}) = \frac{f_{\vett{X} \mid H_1}(\vett{x} \mid H_1)}{f_{\vett{X} \mid H_0}(\vett{x} \mid H_0)}
   \decidi{H_1}{H_0} \eta \qquad\text{con}\qquad P\{L(\vett{X}) > \eta \mid H_0\} = \alpha}
  {\variante{il criterio}{\text{massimizzare } 1 - \beta \text{ con il vincolo } P(\text{errore di prima specie})
     \leq \alpha}
   \variante{dati discreti}{L(\vett{x}) = \frac{p_{\vett{X} \mid H_1}(\vett{x} \mid H_1)}{p_{\vett{X} \mid
     H_0}(\vett{x} \mid H_0)}, \ \text{il più grande falso allarme} \leq \alpha}
   \variante{con il logaritmo, o un'altra funzione crescente}{\Lambda(\vett{x}) = \ln L(\vett{x}) \decidi{H_1}{H_0}
     \ln \eta}
   \variante{nella regola MAP, invece}{\eta = \frac{P(H_0)}{P(H_1)}}}

\voce{qinversa}{Funzione $Q$ inversa}{$Q$ è strettamente decrescente: un solo $x$ per ogni $\alpha \in (0, 1)$}
  {Q^{-1}(\alpha) = x \iff Q(x) = \alpha}
  {\variante{valori}{Q^{-1}(0{,}1) = 1{,}282, \quad Q^{-1}(0{,}05) = 1{,}645, \quad Q^{-1}(0{,}025) = 1{,}960}
   \variante{code più piccole}{Q^{-1}(0{,}01) = 2{,}326, \quad Q^{-1}(0{,}001) = 3{,}090}}

\voce{testnpmedia}{Test di Neyman-Pearson sulla media gaussiana}{$H_0$: $X_i \sim \mathcal{N}(0, \sigma^2)$; \
  $H_1$: $X_i \sim \mathcal{N}(\mu, \sigma^2)$, $\mu > 0$; i.i.d.}
  {\frac{1}{n} \sum_{k=1}^{n} x_k \decidi{H_1}{H_0} \eta' = Q^{-1}(\alpha)\, \frac{\sigma}{\sqrt{n}} \qquad 1 - \beta
   = Q\left( Q^{-1}(\alpha) - \mu\, \frac{\sqrt{n}}{\sigma} \right) \xrightarrow[n \to \infty]{} 1}
  {\variante{la soglia non dipende da $\mu$: il test è il migliore per ogni $\mu > 0$}{}}

\voce{roc}{Curva ROC}{\emph{Receiver Operating Characteristic}: potenza in funzione del falso allarme, al variare della
  soglia}
  {1 - \beta = Q\left( Q^{-1}(\alpha) - d \right), \qquad d = \mu\, \frac{\sqrt{n}}{\sigma}}
  {\variante{$d$ è la distanza fra le due campane, in deviazioni standard della media campione}{}}
